\chapter{下注的目的}

\section{本章导入}

下注，是德州扑克中最重要的主动行动之一。你在牌桌上没有语言可以解释自己的想法，也不能直接要求对手弃牌或支付价值。你真正能使用的工具，只有筹码。你把筹码推进底池，对手必须根据你的下注尺度、行动线、牌面结构和自身牌力作出回应。

从这个角度看，筹码就是你唯一的武器，下注就是你最主要的进攻方式。

很多新手并不是不敢下注，而是不知道自己为什么下注。他们会因为“我有牌”而下注，会因为“我怕对手反超”而下注，会因为“我想知道他有没有”而下注，会因为“我不想让他免费看牌”而下注，也会因为“我觉得这里应该打一下”而下注。表面上看，这些理由都像是下注原因，但它们并不够清晰。

真正有效的下注，必须回到两个核心目的：价值和诈唬。

如果你是价值下注，你希望对手用更差的牌跟注。
如果你是诈唬下注，你希望对手弃掉比你更好的牌。

所有其他概念，例如保护、试探、控池、施压、阻止对手免费看牌，都不能脱离这两个核心目的独立存在。保护性下注本质上仍然要么是在向更差牌和听牌收费，要么是在迫使对手放弃有权益的牌；半诈唬本质上是诈唬的一种，但它同时保留了后续改善成强牌的机会；薄价值下注本质上仍然是价值下注，只是它面对的是更窄、更边缘的更差跟注范围。

本章是整门课程的核心之一。前面我们已经学习了翻前、翻牌、转牌、河牌的处理流程，也学习了不同牌力在范围中的功能。从本章开始，我们进一步追问一个更直接的问题：

\begin{center}
\textbf{我下注到底想让对手做什么？}
\end{center}

本章的核心输出是“下注前必问”：

\begin{center}
\textbf{如果我是价值下注，哪些更差牌会跟？\\
如果我是诈唬，哪些更好牌会弃？\\
如果两个问题都回答不了，就不要轻易下注。}
\end{center}

只要你能在每次下注前问这三个问题，就会立刻减少大量无目的下注、情绪下注和试探性下注。

\section{核心概念}

\subsection{筹码是你唯一的武器}

德州扑克是一项信息不完全的竞技。你不知道对手的底牌，对手也不知道你的底牌。你们之间的信息交流，主要通过行动完成：过牌、下注、跟注、加注、弃牌。所有行动中，下注最具主动性，因为下注会迫使对手立刻作出选择。

当你下注时，对手通常只有几种回应方式：弃牌、跟注、加注。你通过下注改变了底池大小，也改变了对手继续游戏所需要付出的成本。

因此，下注不是简单地“把钱放进底池”。它至少有三层意义：

\begin{enumerate}
\item 它表达了你的范围强度；
\item 它改变了对手继续所需的价格；
\item 它迫使对手用自己的范围作出筛选。
\end{enumerate}

例如，你在翻牌下注 1/3 底池，给对手的压力相对较小，对手可以用较宽范围继续。你在转牌下注 2/3 底池，对手继续成本明显提高，需要放弃更多弱牌。你在河牌使用超池下注，对手往往只能用很强的牌或非常合适的抓诈唬牌继续。

下注的威力来自筹码压力。你下注越大，对手继续范围通常越强；你下注越小，对手可以继续的范围通常越宽。但下注大不一定更好，下注小也不一定更弱。尺度必须服务于你的目的。

新手需要先建立一个基本认识：每一次下注都应该是有目的的筹码投入，而不是凭感觉把筹码推出去。

\subsection{下注的两个核心目的：价值与诈唬}

下注的核心目的主要有两个：价值下注和诈唬下注。

价值下注的目标是让更差的牌跟注。你认为自己的牌领先对手相当一部分继续范围，并且对手会用足够多更差牌支付。此时下注的意义是从这些更差牌中获得更多价值。

例如，你持有 \texttt{AK}，翻牌为 \texttt{A72 rainbow}，面对大盲防守。对手范围中有许多较弱 A、小对子和一些浮动牌。你下注时，希望这些更差牌继续。这就是价值下注。

诈唬下注的目标是让更好的牌弃牌。你当前牌力不足以摊牌获胜，或者摊牌价值很低，但你的下注有机会让对手弃掉比你更好的牌。此时下注的意义是通过弃牌权益赢下底池。

例如，你作为翻前加注者，在 \texttt{A72 rainbow} 牌面持有 \texttt{QJ}。你没有成牌，但这个牌面对你的范围更有利。你小注时，可能让大盲的一些小对子、K 高、Q 高或完全空气牌弃掉。这就是诈唬下注。

价值和诈唬不是由你的主观情绪决定的，而是由对手的继续范围决定的。你说“我是在价值下注”，但如果没有更差牌会跟注，那它就不是好的价值下注。你说“我是在诈唬”，但如果没有更好牌会弃牌，那它就不是好的诈唬。

下注前最重要的问题是：

\begin{center}
\textbf{我希望对手跟注，还是希望对手弃牌？}
\end{center}

如果你希望对手跟注，你必须能说出哪些更差牌会跟。
如果你希望对手弃牌，你必须能说出哪些更好牌会弃。
如果你既说不出更差牌会跟，也说不出更好牌会弃，那么这个下注通常没有意义。

\subsection{价值下注的判断标准}

价值下注看似简单，但真正做好并不容易。很多新手拿到强牌时会下注，但他们很少具体思考对手会用什么更差牌跟注。结果有时下注太小，少拿价值；有时下注太大，把更差牌打跑；有时牌力已经不适合价值下注，却仍然盲目继续。

判断一手牌能否价值下注，可以使用以下标准。

第一，你的牌要领先对手足够多的继续范围。不是说你必须永远领先，而是当对手跟注时，你长期应该能赢过足够比例的跟注牌。

第二，对手范围中要有更差牌愿意跟注。价值来自对手错误或边缘支付。如果对手只会用比你强的牌继续，那么你的下注就是在价值切割自己。

第三，下注尺度要能让目标牌继续。薄价值下注通常不适合过大尺度，因为你的目标是让更差牌舒服地支付；强价值牌则可以根据对手承受能力和底池规划使用更大尺度。

第四，你要考虑被加注后的计划。如果你用一手中等牌力下注，面对加注完全不知道怎么办，那么这次下注可能并不舒服。价值下注不是只考虑对手跟注，也要考虑对手反击。

例如，你持有 \texttt{KQ}，公共牌为 \texttt{K8362}，对手一路过牌。你河牌下注小到中等尺度，希望较弱 K、8x、口袋对子支付，这可能是合理薄价值下注。
但如果公共牌为 \texttt{KQJT9}，你只有一对 K，下注后几乎没有更差牌愿意跟注，而大量顺子、两对、强牌会继续，那么下注就不再是好的价值下注。

价值下注的关键不是“我有牌”，而是“更差牌会不会付钱”。

\subsection{薄价值下注}

薄价值下注是价值下注中的重要能力。它指的是你并不持有非常强的牌，但仍然认为对手会用足够多更差牌跟注，因此下注可以长期盈利。

很多新手在河牌价值不足，就是因为他们只敢用坚果或接近坚果的牌下注。只要牌力不是非常强，他们就选择过牌摊牌。这样虽然减少了一些被更强牌跟注的风险，但也错过了大量更差牌支付的机会。

薄价值下注常见于以下情况：

\begin{itemize}
\item 你持有顶对不错踢脚，对手范围中有较弱顶对；
\item 你持有第二强牌，对手会用第三强牌支付；
\item 对手是跟注站，喜欢用边缘牌看摊牌；
\item 河牌没有明显完成强听牌，牌面较安全；
\item 你下注尺度较小，对手容易用更差牌好奇跟注。
\end{itemize}

薄价值下注的核心不是盲目自信，而是精确估算更差跟注牌。如果对手很紧，只用强牌跟注，那么薄价值下注应该减少；如果对手喜欢跟注，那么薄价值下注应该增加。

例如，你在按钮位 open，大盲跟注。公共牌为 \texttt{K8342}，你持有 \texttt{KQ}，对手河牌过牌。面对很多娱乐玩家，你可以下注小到中等尺度，从 \texttt{KJ}、\texttt{KT}、\texttt{K9}、8x 或部分口袋对子中获取价值。
如果你总是过牌摊牌，虽然经常赢，但会少赢一注。

薄价值下注考验的是牌手对对手跟注范围的判断。它不是为了炫技，而是长期盈利的重要来源。

\subsection{诈唬下注的判断标准}

诈唬下注的目标是让更好的牌弃牌。很多新手理解诈唬时，只关注自己“敢不敢打”，却忽略了对手“会不会弃”。真正的诈唬不是胆量问题，而是条件问题。

一个合理诈唬通常需要满足以下条件。

第一，你的牌摊牌价值较低。如果你的牌本来能通过过牌摊牌赢下不少底池，就不一定适合转成诈唬。适合诈唬的牌通常是错过听牌、空气牌或几乎无法摊牌获胜的牌。

第二，你的行动线能够代表强牌。如果你前面行动一直不像强牌，突然河牌大注，对手很容易不相信。好的诈唬需要故事一致。

第三，牌面支持你的范围。如果公共牌明显更有利于对手，你强行诈唬的可信度下降；如果转牌或河牌强化了你的范围，你的诈唬更容易成功。

第四，对手范围中有足够多可以弃掉的更好牌。比如中等对子、弱顶对、错过但仍领先你的 A 高等。如果对手范围要么是强牌，要么是完全空气牌，诈唬价值就会下降。

第五，对手必须有弃牌能力。对跟注站诈唬，是新手最常见的昂贵错误。对不会弃牌的人，你应该少诈唬，多价值下注。

例如你在翻牌和转牌都代表强范围，河牌完成前门同花，而你手中持有一张关键花色 A，阻断对手坚果同花。对手范围中有大量一对和两对类中等牌力。如果对手有弃牌能力，你就可能拥有一个合理的河牌诈唬机会。

反过来，如果对手是“有一对就看到底”的玩家，那么即使牌面再吓人，诈唬也可能没有意义。

诈唬下注的关键不是“我能不能装强”，而是“对手会不会弃掉比我好的牌”。

\subsection{下注尺度如何影响对手继续范围}

下注尺度会直接影响对手的继续范围。一般来说，下注越小，对手可以用越宽的范围继续；下注越大，对手必须用更强、更有韧性的范围继续。

小注常常会被很多边缘牌跟注，例如小对子、弱对、后门听牌、高牌浮动等。它适合用于高频持续下注、薄价值下注、干燥牌面上的范围下注，以及用较低成本给对手施压。

中注通常开始让对手放弃部分弱牌，同时仍然让一些中等牌力和听牌继续。它适合更明确的价值下注、保护性收费和部分半诈唬。

大注会显著压缩对手继续范围。对手通常需要较强成牌、强听牌或较好的抓诈唬牌才能继续。大注更适合极化策略，也就是用强价值牌和高质量诈唬牌下注。

超池下注则会施加更大压力，常见于一方拥有明显坚果优势，或者希望对对手中等牌力施加最大压力的局面。超池下注不是新手必须频繁使用的工具，但需要理解它的含义：它通常代表非常强的价值牌，或者非常明确的诈唬计划。

不同尺度会让对手继续范围发生变化。比如同一手 \texttt{KQ} 顶对，在干燥牌面上小注，可能被许多较弱牌跟注；如果超池下注，对手继续范围会强得多，你的 \texttt{KQ} 可能就不再适合这样下注。

因此，下注尺度不是随便选的，也不是越大越厉害。尺度必须服务于目的：

\begin{itemize}
\item 想让更差牌广泛跟注，可以用较小尺度；
\item 想向听牌和中等牌收费，可以用中等尺度；
\item 想代表极强牌、压迫对手中等牌力，可以用大尺度；
\item 想最大化坚果优势或制造河牌全下压力，可以考虑超池尺度。
\end{itemize}

下注之前，不仅要问“我下注是价值还是诈唬”，还要问“这个尺度是否服务于我的目的”。

\subsection{小注、中注、大注和超池下注的基本含义}

为了让新手更容易理解，可以先把下注尺度粗略分为四类。

\subsubsection*{小注}

小注通常指 1/4 底池到 1/3 底池左右的下注。它的特点是成本低、频率高、对手继续范围宽。

小注适合：

\begin{itemize}
\item 高牌干燥面上的范围持续下注；
\item 薄价值下注；
\item 让对手用大量较弱牌继续；
\item 用较低成本攻击对手空气牌；
\item 控制底池同时保持主动。
\end{itemize}

小注的缺点是压力有限，较难让中等以上牌力弃牌，也不能充分向强听牌收费。

\subsubsection*{中注}

中注通常指 1/2 底池到 2/3 底池左右的下注。它在价值和压力之间较为平衡。

中注适合：

\begin{itemize}
\item 较强价值牌向更差牌收费；
\item 湿润牌面上让听牌付出代价；
\item 半诈唬时同时保留弃牌权益和后续计划；
\item 转牌继续下注时制造适中压力。
\end{itemize}

中注是实战中最常见的尺度之一。它既不会像小注那样压力太低，也不会像大注那样迅速极化范围。

\subsubsection*{大注}

大注通常指 3/4 底池到接近底池的下注。它会明显压缩对手继续范围，通常更适合强价值牌和强诈唬候选。

大注适合：

\begin{itemize}
\item 你拥有坚果优势；
\item 牌面湿润，需要让强听牌支付较高价格；
\item 你希望对对手中等牌力施加压力；
\item 你有强价值牌并希望构建大底池；
\item 你的诈唬牌有阻断效果或较强后续权益。
\end{itemize}

大注的风险是，如果你的牌力不够强，或诈唬条件不足，很容易让更差牌弃掉、更强牌继续。

\subsubsection*{超池下注}

超池下注是指下注金额超过当前底池。它通常表达非常极化的范围：强到希望打巨大底池的价值牌，以及没有摊牌价值、希望最大化弃牌权益的诈唬牌。

超池下注适合：

\begin{itemize}
\item 你拥有明显坚果优势；
\item 对手范围被限制在大量中等牌力；
\item 你可以代表非常强的牌；
\item 你希望对对手一对、两对等牌力施加极大压力；
\item 你有明确河牌计划或全下计划。
\end{itemize}

初学者不需要频繁使用超池下注，但必须明白：超池下注不是情绪化“吓人”，而是基于范围和坚果优势的高压策略。

\subsection{保护性下注为什么不能脱离价值和诈唬}

新手最常说的一句话是：“我下注是为了保护。”

保护这个概念本身并非错误。很多时候，你的牌当前领先，但对手有听牌或高张可以反超你。下注可以让对手为继续支付价格，避免他免费实现权益。

但保护不能成为一个独立于价值和诈唬之外的模糊理由。你必须进一步问：

\begin{itemize}
\item 我保护的是一手什么牌？
\item 对手有哪些更差牌或听牌会跟注？
\item 我的下注是否让这些牌付出了不合适的价格？
\item 如果对手弃牌，我是否成功让有权益的牌放弃？
\item 如果对手加注，我的牌还能不能继续？
\end{itemize}

例如你持有 \texttt{AK}，翻牌为 \texttt{KJ7 two-tone}。你下注可以让较弱 K、Jx、同花听牌、顺子听牌付费。这种保护性下注本质上同时包含价值：更差牌和听牌会继续，你向它们收费。

但如果你持有 \texttt{K9}，同样牌面下注很大，结果更差牌大多弃牌，更强 K、两对、强听牌继续。你以为自己在保护，实际上可能是在用中等牌力把底池做大。

所以，保护必须落回到两个问题：

\begin{center}
\textbf{更差牌是否会跟注？\\
有权益的更好或接近牌是否会弃牌或付出错误价格？}
\end{center}

如果回答不了，所谓保护很可能只是恐惧驱动的下注。

\subsection{半诈唬}

半诈唬是诈唬的一种，但它和纯诈唬不同。纯诈唬通常依赖对手弃牌；半诈唬即使被跟注，也仍然有机会在后续街改善成强牌。

常见半诈唬包括：

\begin{itemize}
\item 同花听牌下注；
\item 两头顺听牌加注；
\item 组合听牌强力进攻；
\item 带后门同花和顺子可能的高牌持续下注；
\item 一对加听牌的转牌加注。
\end{itemize}

半诈唬有两个赢法：

\begin{enumerate}
\item 对手现在弃牌，你直接赢底池；
\item 对手跟注，你在后续街改善成强牌，赢更大底池。
\end{enumerate}

强听牌适合半诈唬，是因为它们在被跟注时仍然保留较高权益。比如坚果同花听牌加两张高牌，不是完全空气牌。它可以通过下注获得弃牌权益，也可以在被跟注后继续改善。

但半诈唬也不能滥用。弱听牌、低同花听牌、没有位置且面对不弃牌对手时，半诈唬价值会下降。对不会弃牌的人，半诈唬更接近单纯追牌；如果价格不合适，就会亏损。

半诈唬的关键是平衡两个因素：

\begin{center}
\textbf{我有没有足够弃牌权益？\\
如果被跟注，我有没有足够牌面权益？}
\end{center}

两者至少要有一个足够强，半诈唬才有意义。

\subsection{试探性下注为什么是错误概念}

娱乐玩家最常见的下注理由之一，是“我下注是为了试探一下”。

试探性下注的问题在于，它把下注目的变得模糊。你下注后，对手跟注、加注或弃牌，确实会给你一些信息。但这些信息并不是免费的，你为它付出了筹码。而且，对手的回应也不一定能准确告诉你答案。

例如你下注，对手跟注。你可能会想：“他跟了，应该有牌。”但他可能是听牌，可能是弱对，可能是慢打强牌，也可能是浮动。
你下注，对手加注。你可能觉得“他很强”，但他也可能在半诈唬。
你下注，对手弃牌。你赢了底池，但你可能本来就领先，也可能把本来会在后续街诈唬的弱牌打走了。

所以，下注不能只是为了“看看他有什么”。正确的思考应该是：

\begin{itemize}
\item 如果我下注，他哪些更差牌会跟？
\item 如果我下注，他哪些更好牌会弃？
\item 如果他加注，我如何应对？
\item 这个下注是否比过牌更盈利？
\end{itemize}

信息是下注的副产品，不是下注的核心目的。你不能为了买一个模糊信息，就随便投入筹码。

\section{新手常见误区}

\subsection{误区一：为了“知道对手有没有”而下注}

这是最典型的试探性下注。新手拿着中等牌力，不确定自己是否领先，于是下注。他们希望通过对手的反应判断对手有没有强牌。

问题是，对手的反应并不总是诚实而清晰。对手跟注不代表一定强，对手弃牌不代表你原本落后，对手加注也不代表一定是坚果。你用下注换来的信息，经常并不准确。

更重要的是，这类下注往往让你的手牌陷入尴尬：更差牌弃牌，更强牌继续，你既没有价值，也没有成功诈唬。

正确做法是：先判断自己的牌力类别。如果是中等摊牌价值，很多时候过牌控池比试探下注更好。

\subsection{误区二：有牌就下注}

“我有牌，所以我要打。”这是新手最自然的想法，但并不总是正确。

你有底对、第二对子、弱顶对，确实是成牌，但它们是否适合下注，要看更差牌是否会跟注。如果下注后只有更强牌继续，那么你的下注就是错误价值下注。

成牌不等于价值下注。价值下注必须有更差跟注范围。

\subsection{误区三：害怕被反超而过度保护}

很多新手在湿润牌面拿到一对后，会立刻大注，理由是“不让他追”。但如果你的牌只是中等牌力，大注可能会产生相反效果：更差牌被你打跑，更强牌和高权益听牌继续。

保护需要考虑牌力和尺度。强牌可以通过下注向听牌收费；中等牌力则要避免因为恐惧而把底池做得过大。

怕被反超不是下注理由。让对手付出错误价格，才是合理保护的基础。

\subsection{误区四：诈唬不看对手类型}

新手有时会迷恋漂亮的诈唬线，却忽略对手是否会弃牌。对跟注站诈唬，是非常低效的策略。对手如果拿着任何对子都想看摊牌，那么你再合理的故事也很难让他弃牌。

面对不同对手，下注目的要调整：

\begin{itemize}
\item 对跟注站，多价值，少诈唬；
\item 对紧弱玩家，可以增加诈唬，减少薄价值；
\item 对激进玩家，可以用中等摊牌价值诱导；
\item 对被动玩家，面对其突然大注要更谨慎。
\end{itemize}

诈唬不是对所有人都一样有效。下注目的必须与对手类型匹配。

\subsection{误区五：下注尺度随意}

有些新手每次下注都用同一个尺度，不管自己是强牌、弱牌、价值还是诈唬。这会让下注无法服务于目标。

如果你想让较弱牌跟注，尺度过大可能把它们赶走。
如果你想让更好牌弃牌，尺度过小可能没有压力。
如果你想向听牌收费，尺度太小会给对手合适价格。
如果你用中等牌力大注，可能只被更强牌继续。

下注尺度必须和目的相匹配。没有目的的尺度，就是随意下注。

\subsection{误区六：半诈唬和乱诈唬混淆}

半诈唬不是随便拿听牌乱打。它需要有牌面权益，也需要有弃牌权益。

如果你的听牌质量很低，对手又不爱弃牌，你的下注就不是高质量半诈唬，而只是用弱牌扩大底池。反过来，如果你有强听牌、范围优势和位置，半诈唬可以成为非常有力的进攻方式。

半诈唬的关键是：被跟注后仍然有机会，下注时也有机会让对手弃牌。

\section{实战思考框架}

本章的核心框架是“下注前必问”。

\begin{center}
\textbf{如果我是价值下注，哪些更差牌会跟？\\
如果我是诈唬，哪些更好牌会弃？\\
如果两个问题都回答不了，就不要轻易下注。}
\end{center}

\subsection{第一问：如果我是价值下注，哪些更差牌会跟？}

当你准备下注时，首先判断自己是否在价值下注。

如果是价值下注，你必须具体列出更差跟注牌。例如：

\begin{itemize}
\item 我有强顶对，对手会用较弱顶对跟注；
\item 我有两对，对手会用顶对或弱两对跟注；
\item 我有坚果同花，对手会用较小同花或暗三条跟注；
\item 我有薄价值牌，对手会用第二对子或抓诈唬牌跟注。
\end{itemize}

如果你只能说“我可能领先”，但说不出哪些更差牌会跟注，那么下注理由不够充分。

价值下注还要考虑尺度。你要让目标更差牌愿意支付。如果目标牌很弱，你的尺度应该更容易被跟注；如果目标牌较强，或者对手是跟注站，你可以使用更大尺度。

\subsection{第二问：如果我是诈唬，哪些更好牌会弃？}

如果你的下注不是价值下注，那就要判断是否是诈唬。

诈唬下注必须能让更好牌弃牌。你需要具体列出目标弃牌范围。例如：

\begin{itemize}
\item 我希望对手弃掉小对子；
\item 我希望对手弃掉弱顶对；
\item 我希望对手弃掉错过但仍领先我的 A 高；
\item 我希望对手弃掉中等摊牌价值牌。
\end{itemize}

同时，你要判断对手是否真的会弃这些牌。一个理论上应该弃牌但实际上从不弃牌的对手，不是好的诈唬对象。

诈唬还要考虑你的行动线是否可信。你能否代表完成的同花、顺子、两对或强顶对？如果你的故事不一致，对手更容易抓你。

\subsection{第三问：如果两个问题都回答不了，就不要轻易下注}

很多糟糕下注既不是价值，也不是诈唬。

例如你持有中等对子，下注后更差牌大多弃牌，更强牌不会弃牌。这个下注没有价值，也没有诈唬效果。它只是让你投入筹码，却没有让对手犯错。

面对这种局面，过牌通常更好。过牌可以控制底池，保留摊牌价值，观察对手行动，避免把自己放进尴尬位置。

新手要建立一个纪律：

\begin{center}
\textbf{下注之前说不清目的，就先不要下注。}
\end{center}

扑克中不下注也是一种主动选择。过牌不是软弱，弃牌也不是失败。没有目的的下注，才是真正危险的行动。

\subsection{下注目的与牌力类型的对应关系}

不同牌力类型，对应不同下注目的。

\begin{center}
\begin{tabular}{lll}
\toprule
\textbf{牌力类型} & \textbf{常见下注目的} & \textbf{注意事项} \\
\midrule
坚果牌 & 强价值下注 & 构建大底池，避免过度慢打 \\
强价值牌 & 多街价值下注 & 确认更差牌会继续 \\
薄价值牌 & 小到中等价值下注 & 尺度要让更差牌愿意跟 \\
中等摊牌价值 & 多数不宜大注 & 控池，避免价值切割自己 \\
强听牌 & 半诈唬 & 同时依赖弃牌权益和牌面权益 \\
弱听牌 & 谨慎继续 & 不宜用大注硬打 \\
空气牌 & 选择性诈唬 & 必须有弃牌对象和可信故事 \\
\bottomrule
\end{tabular}
\end{center}

这张表可以帮助你把第六章的牌力功能，与本章的下注目的结合起来。

\subsection{下注尺度选择的简化流程}

选择下注尺度时，可以按以下流程思考：

\begin{enumerate}
\item 我是价值还是诈唬？
\item 我的目标牌是强、中等还是弱？
\item 我希望对手继续范围宽一点，还是窄一点？
\item 牌面是干燥还是湿润？
\item 这次下注是否要为下一条街构建底池？
\end{enumerate}

如果你是薄价值，希望较弱牌跟注，通常使用较小尺度。
如果你是强价值，希望构建大底池，可以使用中大尺度。
如果你是半诈唬，希望兼顾弃牌权益和牌面权益，可以根据牌面选择中等或较大尺度。
如果你是纯诈唬，希望让中等牌力弃牌，尺度需要足够有压力。
如果你不知道尺度为什么这样选，说明你还没有明确下注目的。

\section{牌例拆解}

\subsection{牌例一：强价值下注}

\subsubsection*{牌局背景}

有效筹码 100BB。你在按钮位拿到 \texttt{AKo}，前面玩家弃牌，你 open raise 到 2.5BB。大盲跟注。

翻牌为 \texttt{A72 rainbow}。大盲过牌，行动轮到你。

\subsubsection*{新手常见想法}

新手可能会想：

\begin{itemize}
\item “我中了顶对，下注。”
\item “牌面安全，可以打。”
\item “不能让他免费看牌。”
\end{itemize}

这些想法方向没错，但需要更明确下注目的。

\subsubsection*{理性分析}

第一，判断牌力。你持有 \texttt{AK}，在 \texttt{A72 rainbow} 上是强顶对好踢脚，属于强价值牌。

第二，判断下注目的。这是价值下注。你希望大盲用较弱 A、7x、小口袋对子、部分高牌浮动继续。

第三，判断牌面结构。牌面干燥，听牌很少。你不需要为了保护而使用大尺度。小注即可让大量更差牌继续，同时保持范围优势。

第四，选择尺度。1/3 底池左右的小注通常合理。它能让更差牌跟注，也能让对手大量空气牌弃牌。

\subsubsection*{推荐行动}

下注小尺度，进行价值下注。

\subsubsection*{本手牌教训}

强价值下注不是因为“我有牌所以打”，而是因为对手有足够多更差牌会跟注。尺度应服务于让这些牌继续。

\subsection{牌例二：中等牌力不适合试探下注}

\subsubsection*{牌局背景}

有效筹码 100BB。你在 cutoff 拿到 \texttt{JTs}，open raise 到 2.5BB。大盲跟注。

翻牌为 \texttt{KJ7 two-tone}。大盲过牌，行动轮到你。

\subsubsection*{新手常见想法}

很多新手会想：

\begin{itemize}
\item “我有第二对子，下注看看他有没有 K。”
\item “牌面有听牌，不能给他免费看。”
\item “如果他跟注，我再判断。”
\end{itemize}

这就是典型的试探性下注。

\subsubsection*{理性分析}

第一，判断牌力。你持有第二对子，有一定摊牌价值，但不是强价值牌。

第二，判断价值下注。下注后，哪些更差牌会跟注？可能有一些 7x、听牌或较弱 Jx，但很多更差空气牌会弃牌。更强 Kx、强 J、两对、暗三条和强听牌会继续。

第三，判断诈唬下注。你下注能让哪些更好牌弃牌？大多数 Kx 不会因为一枪就弃牌。更强 J 也未必弃。你的下注很难让足够多更好牌弃牌。

第四，比较过牌。过牌可以控制底池，保留摊牌价值，也让对手的空气牌和错过听牌留在范围里。

\subsubsection*{推荐行动}

倾向过牌控池。

如果对手非常被动且会用大量更差牌跟注，可以考虑小注；但默认情况下，不应为了“试探”而下注。

\subsubsection*{本手牌教训}

中等摊牌价值最容易被试探性下注误用。下注前必须问：更差牌会跟吗？更好牌会弃吗？如果都不明确，过牌通常更好。

\subsection{牌例三：半诈唬下注}

\subsubsection*{牌局背景}

有效筹码 100BB。你在按钮位拿到 \texttt{AQs}，前面玩家弃牌，你 open raise 到 2.5BB。大盲跟注。

翻牌为 \texttt{J T 4}，其中两张与你同花。大盲过牌，行动轮到你。

\subsubsection*{新手常见想法}

一些新手会想：

\begin{itemize}
\item “我还没中牌，等中了再打。”
\item “如果下注被跟，我现在还是 A 高。”
\item “听牌就跟，不要主动打。”
\end{itemize}

这些想法低估了半诈唬的价值。

\subsubsection*{理性分析}

第一，判断牌力。你有坚果同花听牌、两张高牌和顺子相关可能，是强听牌。

第二，判断下注目的。这不是纯价值下注，而是半诈唬。你希望对手弃掉部分小对子、弱 4x、没有强权益的空气牌；如果对手跟注，你仍然有大量转牌和河牌改善机会。

第三，判断牌面。牌面较湿润，对手有继续范围，但你的牌权益很高，适合主动进攻。

第四，选择尺度。可以使用中等尺度，既制造弃牌权益，也为后续完成强牌构建底池。

\subsubsection*{推荐行动}

下注半诈唬。

\subsubsection*{本手牌教训}

半诈唬不是乱 bluff。强听牌下注有两个赢法：对手现在弃牌，或者你后续改善成强牌。

\subsection{牌例四：薄价值下注}

\subsubsection*{牌局背景}

有效筹码 100BB。你在按钮位拿到 \texttt{KQ}，open raise 到 2.5BB。大盲跟注。

翻牌为 \texttt{K83 rainbow}。大盲过牌，你下注 1/3 底池，大盲跟注。

转牌为 \texttt{6}。双方过牌。

河牌为 \texttt{2}。大盲过牌，行动轮到你。

\subsubsection*{新手常见想法}

新手可能会想：

\begin{itemize}
\item “我有顶对，过牌也能赢。”
\item “如果他有两对或暗三条，我下注会输钱。”
\item “河牌没必要冒险。”
\end{itemize}

这类想法会导致漏掉薄价值。

\subsubsection*{理性分析}

第一，判断牌力。你持有 \texttt{KQ}，在 \texttt{K8362} 上是较强顶对。转牌双方过牌后，对手范围并不一定很强。

第二，判断价值下注。更差牌包括较弱 K、8x、部分口袋对子，甚至一些好奇跟注的 A 高。面对小到中等尺度，这些牌可能会支付。

第三，判断尺度。因为这是薄价值下注，尺度不宜过大。小到中等下注更容易被更差牌跟注。

第四，判断风险。你会输给少量两对和暗三条，但这不代表不能下注。价值下注不是要求永远领先，而是要求被跟注时长期领先足够多。

\subsubsection*{推荐行动}

下注小到中等尺度，进行薄价值下注。

\subsubsection*{本手牌教训}

过牌能赢，不代表过牌最好。如果更差牌会支付，河牌应该争取薄价值。

\subsection{牌例五：错误的河牌诈唬对象}

\subsubsection*{牌局背景}

有效筹码 100BB。你在按钮位拿到 \texttt{QTs}，open raise 到 2.5BB。大盲是一名明显跟注站，喜欢用任何对子跟到河牌。

翻牌为 \texttt{K72 rainbow}。大盲过牌，你下注 1/3 底池，大盲跟注。

转牌为 \texttt{4}。大盲过牌，你继续下注 2/3 底池，大盲跟注。

河牌为 \texttt{9}。你没有成牌。大盲过牌，行动轮到你。

\subsubsection*{新手常见想法}

新手可能会想：

\begin{itemize}
\item “我已经打了两枪，第三枪才能把故事讲完。”
\item “他可能只有 7 或小对子，河牌大注能打走。”
\item “我不打就一定输了。”
\end{itemize}

这些想法忽略了对手类型。

\subsubsection*{理性分析}

第一，判断摊牌价值。你没有成牌，几乎没有摊牌价值。

第二，判断诈唬目标。你希望对手弃掉 7x、小对子、弱 K 或其他比你好的牌。

第三，判断对手是否会弃。关键问题是：这个大盲是跟注站。他已经用很多边缘牌跟到河牌，且倾向于用任何对子看摊牌。对这样的对手，诈唬成功率很低。

第四，判断行动线。虽然你可以代表部分 Kx 和强牌，但对手如果不关心你的故事，诈唬就没有意义。

\subsubsection*{推荐行动}

过牌放弃。

\subsubsection*{本手牌教训}

诈唬必须选择会弃牌的对象。对不会弃牌的人，少诈唬，多价值下注。好的故事讲给听得懂的人才有意义。

\section{本章总结}

本章的核心，是让你真正理解下注不是凭感觉把筹码推出去，而是带有明确目的的进攻行动。

筹码是你唯一的武器。你通过下注表达范围、改变价格、迫使对手筛选继续范围。但筹码不能随意使用，每一次下注都应该服务于具体目标。

下注的两个核心目的，是价值和诈唬。

价值下注的目标，是让更差牌跟注。
诈唬下注的目标，是让更好牌弃牌。

保护、半诈唬、薄价值、施压、阻止免费看牌等概念，都不能脱离这两个核心目的独立存在。保护本质上要么是让更差牌和听牌付费，要么是迫使有权益的牌放弃；半诈唬是带有牌面权益的诈唬；薄价值是边缘但仍然盈利的价值下注。

下注尺度必须服务于下注目的。小注让对手继续范围更宽，适合薄价值、高频下注和低成本施压；中注平衡价值与压力；大注压缩对手继续范围，更适合强价值和极化诈唬；超池下注则通常需要坚果优势和明确的极化策略。

本章的核心框架是：

\begin{center}
\textbf{如果我是价值下注，哪些更差牌会跟？\\
如果我是诈唬，哪些更好牌会弃？\\
如果两个问题都回答不了，就不要轻易下注。}
\end{center}

新手最需要戒掉的是试探性下注。下注不是为了“看看对手有没有”，因为对手的反应并不总是清晰可靠。信息只是下注的副产品，不是下注的目的。

下一章我们将讨论“激进的游戏”。当你理解下注目的之后，就要进一步理解主动权、弃牌权益和进攻频率。德州扑克不是只靠摊牌赢底池的游戏，优秀牌手必须学会在合适条件下主动争取底池。

\section{课后训练}

\subsection*{训练一：下注目的记录}

下一次打牌后，记录 10 手你主动下注的牌。每一手都回答：

\begin{enumerate}
\item 我下注是价值还是诈唬？
\item 如果是价值下注，哪些更差牌会跟注？
\item 如果是诈唬下注，哪些更好牌会弃牌？
\item 我的下注尺度是否服务于这个目的？
\item 如果对手加注，我当时有没有计划？
\end{enumerate}

如果你无法回答第 2 或第 3 个问题，说明这次下注目的不清。

\subsection*{训练二：找出试探性下注}

从最近牌局中找出 5 手你“为了看看对手有没有”而下注的牌。

逐手分析：

\begin{itemize}
\item 我的牌属于什么牌力类型？
\item 更差牌是否会跟注？
\item 更好牌是否会弃牌？
\item 过牌是否更合理？
\item 下注后我获得的信息是否真的可靠？
\end{itemize}

这个训练的目标，是逐步戒掉试探性下注。

\subsection*{训练三：薄价值下注训练}

找出 10 手你河牌过牌摊牌并赢下底池的牌，判断其中哪些可以薄价值下注。

重点回答：

\begin{enumerate}
\item 对手是否可能用更差牌跟注？
\item 我应该使用小注、中注还是大注？
\item 对手是跟注站、紧弱玩家，还是正常玩家？
\item 我当时过牌是合理控池，还是害怕下注？
\end{enumerate}

训练目标是提高河牌价值获取能力。

\subsection*{训练四：诈唬对象选择}

记录 5 手你尝试诈唬的牌，逐手回答：

\begin{itemize}
\item 对手是否有弃牌能力？
\item 我的行动线是否能代表强牌？
\item 牌面是否支持我的范围？
\item 我的牌是否几乎没有摊牌价值？
\item 我的下注尺度是否足够让更好牌弃牌？
\end{itemize}

如果对手不具备弃牌能力，未来应减少类似诈唬。

\subsection*{训练五：下注尺度复盘}

找出最近 10 次下注，记录你的下注尺度，并判断它是否符合目的：

\begin{itemize}
\item 小注：是否为了让更宽范围继续或低成本施压？
\item 中注：是否为了价值和保护之间平衡？
\item 大注：是否因为我有强价值或强诈唬候选？
\item 超池下注：是否真的有坚果优势和极化理由？
\end{itemize}

这个训练的目标，是让下注尺度从习惯动作变成策略选择。

\subsection*{训练六：半诈唬判断}

记录 10 手你持有听牌并主动下注或加注的牌，分析它们是否是真正的半诈唬。

逐手回答：

\begin{enumerate}
\item 我是否有足够牌面权益？
\item 对手是否有足够弃牌率？
\item 我的听牌是强听牌还是弱听牌？
\item 如果被跟注，转牌或河牌有哪些好牌？
\item 如果对手加注，我是否有继续计划？
\end{enumerate}

不要把所有听牌下注都称为半诈唬。只有同时具备弃牌权益和牌面权益的下注，才是高质量半诈唬。
